%
%
\section{In-Context Reinforcement Learning}

% \red{In this work, we focus on ICRL under random policies and random contexts.}\blue{I would remove this.} 
%
This section introduces the background of ICRL mechanisms and three SOTA ICRL algorithms. We start by presenting the preliminaries of ICRL.

\subsection{Preliminaries}

RL problems are generally formulated as Markov Decision Processes (MDPs)~\citep{sutton2018reinforcement}.
% \blue{I understand that we want to push our paper and get citations, but I would remove it from here. This is a generic citation for RL, I don't think other people would cite our work in that context. If we were discussing safe RL I wouldn't have a problem with this.}
An MDP can be represented by a tuple $\tau = (S, A, R, P, \rho)$, where $S$ and $A$ denote finite state and action spaces, $R: S \times A \to \Real$ denotes the reward function that evaluates the quality of the action, $P: S \times A \times S \to [0, 1]$ denotes the transition probability that describes the dynamics of the system, and $\rho: S \to [0, 1]$ denotes the initial state distribution.  



A policy $\pi$ defines a mapping from states to probability distributions over actions, providing a strategy that guides the agent in the decision making.
%
The agent interacts with the environment following the policy $\pi$ and the transition dynamics of the system, and then generates an episode of the transition data $(s_0, a_0, r_0, \cdots)$.
%
The performance measure $J (\pi)$ is defined by the expected discounted cumulative reward under the policy $\pi$
\begin{align}
    J (\pi) = \mbE_{s_0 \sim \rho, a_t \sim \pi(\cdot | s_t), s_{t+1} \sim P(\cdot|s_t, a_t)} \left[ \sum_{t=0}^\infty \gamma^t r_t \right].
\end{align}
%
The goal of RL is to find an optimal policy $\pi^*$ that maximizes $J (\pi)$.
%
It is crucial to recognize that $\pi^*$ often varies across different MDPs (environments). Thus, the optimal policy for standard RL must be re-learned each time a new environment is encountered.
%
%
Under this circumstance, ICRL proposes to pretrain a FM on a wide variety of pretraining environments, and then deploy it in the \textit{unseen} test environments without updating parameters in the pretrained model, i.e., zero-shot generalization~\citep{sohn2018hierarchical, mazoure2022improving, zisselman2023explore, kirk2023survey}.



%%%%%%%%%%%%%%%%%%%%%%%%%%%%%%%%%%%%%%%%%%%%%%%%%%%%%%%%%%%%%%%%%%%%%%%%

\subsection{Supervised Pretraining Mechanism}

In this subsection, we introduce the methodology behind ICRL--a supervised pretraining mechanism. Consider two distributions over environments $\ccalT_{\text{train}}$ and $\ccalT_{\text{test}}$ {for} pretraining and test (evaluation), respectively. Each environment, along with its corresponding MDP $\tau$, can be regarded as an instance drawn from the environment distributions, where each environment may exhibit distinct reward functions and transition dynamics.
%
Given an environment $\tau$, a context/prompt $\ccalC = \{s_i, a_i, r_i, s_i'\}_{i \in [n]}$ refers to a sample from a pretraining context dataset $D_{\text{train}} (\cdot \mid \tau)$, i.e., $\ccalC \sim D_{\text{train}} (\cdot \mid \tau)$, which are collected through interactions between a behavior policy and the environment $\tau$ (see e.g., Algorithm~\ref{alg_collect_context}). Notably, $D_{\text{train}} (\cdot \mid \tau)$ contains the contextual information regarding the environment $\tau$.
%
We next consider a query state distribution $D_q^\tau$ and a label policy that maps the query state to the distribution of the action label, i.e., $\pi_{l}: S \to \Delta_{a_l} (A)$, where $\Delta_{a_l}(\cdot)$ denotes the probability distribution and $a_l$ denotes the prediction/output goal of the FM.
 %\blue{Is the output a probability or an action?} \orange{$\Delta$ generally represents the distribution in the literature.} \blue{But the label should be an action, not a probability.}
Then, the joint distribution over the environment $\tau$, context $\ccalC$, query state $s_{q}$, and action label $a_{l}$ is given by
\begin{align}\label{eqn_pretrain_distribution}
    P_{\text{train}} (\tau, \ccalC, s_{q}, a_{l}) = \ccalT_{\text{train}} (\tau) \cdot D_{\text{train}} (\ccalC | \tau) \cdot D_q^\tau \cdot \pi_{l}(a_{l} | s_{q}).
\end{align}
%
%
%
ICRL follows a supervised pretraining mechanism. More concretely, a FM with parameter $\theta$ (denoted by $\ccalF_\theta: \mbC \times S \to \Delta (a_{l})$) is pretrained to predict the action label $a_{l}$ given the context $\ccalC$ and query state $s_{q}$. To do so, the current literature~\citep{laskin2022context, lee2024supervised, dong2024context} often considers the following objective function
%
\begin{align}\label{eqn_loss}
    \theta^* = \argmin_\theta \mbE_{P_{\text{train}}} \left[ l \left( \ccalF_\theta(\cdot \mid \ccalC, s_{q}), a_{l}\right) \right],
\end{align}
%
where $l(\cdot, \cdot)$ represents the loss function, for instance, negative log-likelihood (NLL) for discrete-action problems and mean square error (MSE) for continuous-action scenarios.
%\blue{You need explain when you mention the FM, what is $\mathcal{F}_\theta$}
%
We note that while the current SOTA ICRL algorithms (AD, DPT, and DIT) adhere to a common objective function \eqref{eqn_loss}, they differ significantly in constructing the context, query state, and action label.
%
Furthermore, it is important to highlight that selecting appropriate action labels in existing ICRL algorithms can be prohibitively expensive. We discuss this in more detail for each of the ICRL baselines  that follow.





\paragraph{Algorithm Distillation.} Instead of learning an optimal policy for a specific environment, AD proposes to learn a RL algorithm itself across a wide range of environments. This is, the FM in AD is pretrained to imitate the source (or behavior) algorithm over the pretraining environment distribution $\ccalT_{\text{train}}$. In general, AD demands a well-trained source algorithm with its complete learning history (from the initial policy to the final-trained policy). Additionally, AD is restricted to the environments with short episode length, as the context must capture cross-episode information of the source algorithm, while standard transformers often have limited context length.
%
In terms of the objective function \eqref{eqn_loss}, AD takes the state at the time step $t$ as the query state, $a_t$ from the source algorithm as the action label, and the episodic history data $(s_0, a_0, r_0, \cdots, s_{t-1}, a_{t-1}, r_{t-1})$ to be the context.
%





\paragraph{Decision Pretrained Transformer.} 
%
Instead of being stringent on the context and the environment itself, DPT handles the context and query state in a more general manner.
%
Specifically, in terms of the objective function \eqref{eqn_loss}, DPT can consider a random collection of transitions as the context $\ccalC$, a random query state $s_q$ drawn from $D_q$, and an optimal action label corresponding to $s_q$.
%
Despite less requirements on the context, DPT necessitates access to optimal policies for the optimal action labels in all pretraining environments, which may not be available in real-world applications.






\paragraph{Decision Importance Transformer.}

DIT proposes to learn ICRL without optimal action labels, following the same supervised pretraining mechanism as DPT.
%
To that end, DIT chooses to consider every possible state and action in the context as the query state and corresponding action label. 
%
It is important to point out that DIT requires a (partially) complete episode to form the context, enabling the computation of a return-to-go for each state-action pair. By mapping the return-to-go to a weight that reflects the quality of each state-action pair, DIT can pretrain the FM using the DPT structure, augmented by the weight assigned to each query state and action label. 
%
In other words, DIT prioritizes the training on high-return pairs.
%
Notably, DIT still mandates that more than $30\%$ of the context data comes from well-trained policies to ensure the coverage of good action labels.



% \red{However, both AD and DPT rely on optimal (or well-trained) source policies. When only suboptimal policies are available, these SOTA algorithms encounter significant difficulties in achieving effective ICRL. ~\blue{[REF]} proposes Decision Importance Transformer (DIT) that enables to learn ICRL with suboptimal policies.}
% \vspace{10pt}

To summarize, it is worth highlighting that all these SOTA ICRL algorithms necessitate varying degrees of well-trained, or even optimal policies during the pretraining phase.
%
However, obtaining such policies for real-world applications is often prohibitive, as it demands extensive training across a vast number of real-world environments.
%
This challenge becomes even more pronounced in the domains where the transition data exhibits high variance to train an effective policy.
%
% It is then expected for these ICRL algorithm to encounter significant difficulties when only suboptimal (or even random) policies and random contexts are available.
%
Thus, executing ICRL under (e.g., uniform) random policies and random contexts is crucial for enabling the practical application of ICRL in the real-world.






%%%%%%%%%%%%%%%%%%%%%%%%%%%%%%%%%%%%%%%%%%%%%%%%%%%%%%%%%%%%%%%%%%%%%%%%
